\documentclass[10pt,a4paper]{amsart}
\usepackage{amsmath,amssymb,amsthm,graphicx}
% Retained preamble for the Stacks Project chapter snapshot in this directory.
% The chapter's own preamble.tex is a copy of the shared Stacks Project preamble; the block
% below is sliced MECHANICALLY, line for line, from it. Nothing was added, removed or
% redefined; the wrapper class, the xr-hyper/hyperref wiring and the externaldocument list
% are omitted because the assembler supplies the document class.
%
% lines taken from preamble.tex: 3, 5, 9, 11, 12, 156, 157, 158, ... (51 lines total)
%
% package coverage restored from the chapter's own preamble (2026-10-13): xy, xy-2cell

\usepackage[all]{xy}
\xyoption{2cell}
\UseAllTwocells

\newenvironment{reference}{\comment}{\endcomment}
\newenvironment{slogan}{\comment}{\endcomment}
\newenvironment{history}{\comment}{\endcomment}
\theoremstyle{plain}
\newtheorem{theorem}[subsection]{Theorem}
\newtheorem{proposition}[subsection]{Proposition}
\newtheorem{lemma}[subsection]{Lemma}
\theoremstyle{definition}
\newtheorem{definition}[subsection]{Definition}
\newtheorem{example}[subsection]{Example}
\newtheorem{exercise}[subsection]{Exercise}
\newtheorem{situation}[subsection]{Situation}
\theoremstyle{remark}
\newtheorem{remark}[subsection]{Remark}
\newtheorem{remarks}[subsection]{Remarks}
\numberwithin{equation}{subsection}
\def\lim{\mathop{\mathrm{lim}}\nolimits}
\def\colim{\mathop{\mathrm{colim}}\nolimits}
\def\Spec{\mathop{\mathrm{Spec}}}
\def\Hom{\mathop{\mathrm{Hom}}\nolimits}
\def\Ext{\mathop{\mathrm{Ext}}\nolimits}
\def\SheafHom{\mathop{\mathcal{H}\!\mathit{om}}\nolimits}
\def\SheafExt{\mathop{\mathcal{E}\!\mathit{xt}}\nolimits}
\def\Sch{\mathit{Sch}}
\def\Mor{\mathop{\mathrm{Mor}}\nolimits}
\def\Ob{\mathop{\mathrm{Ob}}\nolimits}
\def\Sh{\mathop{\mathit{Sh}}\nolimits}
\def\NL{\mathop{N\!L}\nolimits}
\def\CH{\mathop{\mathrm{CH}}\nolimits}
\def\proetale{{pro\text{-}\acute{e}tale}}
\def\etale{{\acute{e}tale}}
\def\QCoh{\mathit{QCoh}}
\def\Ker{\mathop{\mathrm{Ker}}}
\def\Im{\mathop{\mathrm{Im}}}
\def\Coker{\mathop{\mathrm{Coker}}}
\def\Coim{\mathop{\mathrm{Coim}}}
\DeclareMathSymbol{\boxtimes}{\mathbin}{AMSa}{"02}
\def\QCohstack{\mathcal{QC}\!\mathit{oh}}
\def\Cohstack{\mathcal{C}\!\mathit{oh}}
\def\Spacesstack{\mathcal{S}\!\mathit{paces}}
\def\Quotfunctor{\mathrm{Quot}}
\def\Hilbfunctor{\mathrm{Hilb}}
\def\Curvesstack{\mathcal{C}\!\mathit{urves}}
\def\Polarizedstack{\mathcal{P}\!\mathit{olarized}}
\def\Complexesstack{\mathcal{C}\!\mathit{omplexes}}
\def\Pic{\mathop{\mathrm{Pic}}\nolimits}
\def\Picardstack{\mathcal{P}\!\mathit{ic}}
\def\Picardfunctor{\mathrm{Pic}}
\def\Deformationcategory{\mathcal{D}\!\mathit{ef}}
\usepackage{cleveref}
\crefname{originalproposition}{original proposition}{original propositions}
\Crefname{originalproposition}{Original proposition}{Original propositions}
\makeatletter
\newlabel{constructions-section-relative-glueing}{{01LG}{source}{}{sourcepointer-1}{}}
\newlabel{constructions-section-relative-glueing@cref}{{[originalproposition][0][]01LG}{[1][0][]source}{}{}{}}
\makeatother
\crefname{originalproposition}{original proposition}{original propositions}
\Crefname{originalproposition}{Original proposition}{Original propositions}
\makeatletter
\newlabel{lemma-hilb-d-An}{{0B99}{source}{}{sourcepointer-2}{}}
\newlabel{lemma-hilb-d-An@cref}{{[originalproposition][0][]0B99}{[1][0][]source}{}{}{}}
\makeatother
\crefname{originalproposition}{original proposition}{original propositions}
\Crefname{originalproposition}{Original proposition}{Original propositions}
\makeatletter
\newlabel{lemma-hilb-d-sheaf}{{0B95}{source}{}{sourcepointer-3}{}}
\newlabel{lemma-hilb-d-sheaf@cref}{{[originalproposition][0][]0B95}{[1][0][]source}{}{}{}}
\makeatother
\crefname{originalproposition}{original proposition}{original propositions}
\Crefname{originalproposition}{Original proposition}{Original propositions}
\makeatletter
\newlabel{morphisms-lemma-image-proper-scheme-closed}{{01W6}{source}{}{sourcepointer-4}{}}
\newlabel{morphisms-lemma-image-proper-scheme-closed@cref}{{[originalproposition][0][]01W6}{[1][0][]source}{}{}{}}
\makeatother
\crefname{originalproposition}{original proposition}{original propositions}
\Crefname{originalproposition}{Original proposition}{Original propositions}
\makeatletter
\newlabel{schemes-lemma-glue-functors}{{01JJ}{source}{}{sourcepointer-5}{}}
\newlabel{schemes-lemma-glue-functors@cref}{{[originalproposition][0][]01JJ}{[1][0][]source}{}{}{}}
\makeatother
\crefname{originalproposition}{original proposition}{original propositions}
\Crefname{originalproposition}{Original proposition}{Original propositions}
\makeatletter
\newlabel{schemes-lemma-immersion-when-closed}{{01IQ}{source}{}{sourcepointer-6}{}}
\newlabel{schemes-lemma-immersion-when-closed@cref}{{[originalproposition][0][]01IQ}{[1][0][]source}{}{}{}}
\makeatother
\title{Hilb\textasciicircum{}d of a scheme with affine-opens through points is representable}
\author{The Stacks Project Authors}
\date{stacks-project:pic.tex:local-snapshot-20261012; GNU FDL 1.2 or later}
\begin{document}
\maketitle
\noindent\textit{Editorial note.} Statement span is lines 515-522 of sources/stacks-pic/pinned-original-source.tex; proof span is lines 524-586, the complete author proof (61 body lines).\par\medskip
\noindent\textit{Editorial note.} Navigation stubs: 6 cited-but-not-included labels are supplied as non-load-bearing pointers with their pinned tags where the snapshot tag index has them.\par\medskip
\noindent\textit{Editorial note.} extra\_packages is empty on purpose: the retained preamble already loads xy with options and a second declaration clashes in the assembler wrapper.\par\medskip
\noindent\textit{Editorial note.} Text is retained verbatim, including the author's own typos and placeholders.\par\medskip
\noindent\textit{Editorial note.} Queue index 4 of queue\_pic.json (pb=3101, stmt 515-522, proof 524-586).\par\medskip
\noindent\textit{Editorial note.} stacks-pic wave-1 (pi01): fresh adjudication accept (2026-10-14); pool per standing queue hypercovering -> pic (SMC reassigned to Zhen per Chief 08:5x); pic.tex 27 labels/23 with proof judged 21 accept / 2 reject / 4 pb=0 env-dup rows excluded; spans cite-free and reference/slogan/history-free so no patchenv needed. Named omission: Author omits the details of the reduction to affine S (two standard alternative routes named) - a named routine reduction; the F\_U gluing argument including the non-separated footnote is in-span.\par\medskip
\noindent\textit{Allowed standard prerequisites (editorial).}\begin{itemize}
\item Explicit prerequisite boundary. The cited inputs are prerequisites; the argument in the retained span is not replaced by any cited result.
\end{itemize}\medskip
\ifdefined\hypertarget\hypertarget{sourcepointer-1}{}\fi
\noindent\textit{Source reference.} Original proposition 01LG: sources/stacks-pic/pinned-original-source.tex, source lines 586--586. Cited by the retained statement or proof but not part of the collected spans; supplied only as a navigation pointer. The argument does not depend on its content.\par\medskip
\ifdefined\hypertarget\hypertarget{sourcepointer-2}{}\fi
\noindent\textit{Source reference.} Original proposition 0B99: sources/stacks-pic/pinned-original-source.tex, source lines 411--414. Cited by the retained statement or proof but not part of the collected spans; supplied only as a navigation pointer. The argument does not depend on its content.\par\medskip
\ifdefined\hypertarget\hypertarget{sourcepointer-3}{}\fi
\noindent\textit{Source reference.} Original proposition 0B95: sources/stacks-pic/pinned-original-source.tex, source lines 282--286. Cited by the retained statement or proof but not part of the collected spans; supplied only as a navigation pointer. The argument does not depend on its content.\par\medskip
\ifdefined\hypertarget\hypertarget{sourcepointer-4}{}\fi
\noindent\textit{Source reference.} Original proposition 01W6: sources/stacks-pic/pinned-original-source.tex, source lines 586--586. Cited by the retained statement or proof but not part of the collected spans; supplied only as a navigation pointer. The argument does not depend on its content.\par\medskip
\ifdefined\hypertarget\hypertarget{sourcepointer-5}{}\fi
\noindent\textit{Source reference.} Original proposition 01JJ: sources/stacks-pic/pinned-original-source.tex, source lines 586--586. Cited by the retained statement or proof but not part of the collected spans; supplied only as a navigation pointer. The argument does not depend on its content.\par\medskip
\ifdefined\hypertarget\hypertarget{sourcepointer-6}{}\fi
\noindent\textit{Source reference.} Original proposition 01IQ: sources/stacks-pic/pinned-original-source.tex, source lines 586--586. Cited by the retained statement or proof but not part of the collected spans; supplied only as a navigation pointer. The argument does not depend on its content.\par\medskip
\begin{proposition}
\label{proposition-hilb-d-representable}
Let $X \to S$ be a morphism of schemes. Let $d \geq 0$. Assume
for all $(s, x_1, \ldots, x_d)$ where $s \in S$ and
$x_1, \ldots, x_d \in X_s$ there exists an affine open $U \subset X$
with $x_1, \ldots, x_d \in U$. Then $\Hilbfunctor^d_{X/S}$ is
representable by a scheme.
\end{proposition}

\begin{proof}
Either using relative glueing (Constructions, Section
\ref{constructions-section-relative-glueing}) or using
the functorial point of view
(Schemes, Lemma \ref{schemes-lemma-glue-functors})
we reduce to the case where $S$ is affine. Details omitted.

\medskip\noindent
Assume $S$ is affine. For $U \subset X$ affine open, denote
$F_U \subset \Hilbfunctor^d_{X/S}$ the subfunctor such that
for a scheme $T/S$ an element $Z \in \Hilbfunctor^d_{X/S}(T)$
is in $F_U(T)$ if and only if $Z \subset U_T$. We will use
Schemes, Lemma \ref{schemes-lemma-glue-functors}
and the subfunctors $F_U$ to conclude.

\medskip\noindent
Condition (1) is Lemma \ref{lemma-hilb-d-sheaf}.

\medskip\noindent
Condition (2)(a) follows from the fact that $F_U = \Hilbfunctor^d_{U/S}$
and that this is representable by Lemma \ref{lemma-hilb-d-An}.
Namely, if $Z \in F_U(T)$, then $Z$ can be viewed as a closed subscheme
of $U_T$ which is finite locally free of degree $d$ over $T$ and hence
$Z \in \Hilbfunctor^d_{U/S}(T)$. Conversely, if $Z \in \Hilbfunctor^d_{U/S}(T)$
then $Z \to U_T \to X_T$ is a closed immersion\footnote{This is clear
if $X \to S$ is separated as in this case Morphisms, Lemma
\ref{morphisms-lemma-image-proper-scheme-closed}
tells us that the immersion $\varphi : Z \to X_T$ has closed image
and hence is a closed immersion by
Schemes, Lemma \ref{schemes-lemma-immersion-when-closed}. We suggest the
reader skip the rest of this footnote as we don't know of any instance
where the assumptions on $X \to S$ hold but $X \to S$ is not separated.
In the general case, let $x \in X_T$ be a point in the closure of
$\varphi(Z)$. We have to show that $x \in \varphi(Z)$. Let $t \in T$ be the
image of $x$. By assumption on $X \to S$ we can choose an affine open
$W \subset X_T$ containing $x$ and $\varphi(Z_t)$. Then $\varphi^{-1}(W)$
is an open containing the whole fibre $Z_t$ and since $Z \to T$ is closed,
we may after replacing $T$ by an open neighbourhood of $t$ assume that
$Z = \varphi^{-1}(W)$. Then $\varphi(Z) \subset W$ is closed by the
separated case (as $W \to T$ is separated) and we conclude $x \in \varphi(Z)$.}
and we may view $Z$ as an element of $F_U(T)$.

\medskip\noindent
Let $Z \in \Hilbfunctor^d_{X/S}(T)$ for some scheme $T$ over $S$. Let
$$
B = (Z \to T)\left((Z \to X_T \to X)^{-1}(X \setminus U)\right)
$$
This is a closed subset of $T$ and it is clear that over the open
$T_{Z, U} = T \setminus B$ the restriction $Z_{t'}$ maps into $U_{T'}$.
On the other hand, for any $b \in B$ the fibre $Z_b$ does not map
into $U$. Thus we see that given a morphism $T' \to T$ we
have $Z_{T'} \in F_U(T')$ $\Leftrightarrow$ $T' \to T$ factors through
the open $T_{Z, U}$. This proves condition (2)(b).

\medskip\noindent
Condition (2)(c) follows from our assumption on $X/S$. All we have
to do is show the following: If $T$ is the spectrum of a field
and $Z \subset X_T$ is a closed subscheme, finite flat of degree
$d$ over $T$, then $Z \to X_T \to X$ factors through an affine open
$U$ of $X$. This is clear because $Z$ will have at most $d$ points
and these will all map into the fibre of $X$ over the image point
of $T \to S$.
\end{proof}

\end{document}
